\chapter{Fundamentals}
\nopagecolor{}
\label{chap:chap1}
\vspace{-50pt}%

This chapter presents the fundamental characteristics of the Ymir programming language. It introduces the first necessary steps to build a Ymir program, and is also a first step into the basics of low level programming via a high level programming language. This chapter also presents the basics of how this book is structured, and guides the reader towards a full understanding of the concepts it presents.

\minitoc%

The development of a program using the Ymir language follows three fundamental
steps. First, a text editor is used to create text files named source files, in
which source code following the Ymir grammar and semantics is written. These
source files are passed to the compiler, which transforms source code
information into object files made up of machine code. These object files are
finally linked by a linker to form an executable file. The linker also pulls in
the code of the libraries that the program uses, such as the standard library,
which are distributed as archives of already compiled object files. In practice,
a single invocation of the compiler performs both the compilation and the
linking steps.

\input{chapters/chapter1/figures/compilation_chain}
\vfill%
\pagebreak%
          
A source file is a text file encoded in \textit{utf-8} with the
extension \textit{.yr}. It is important to use the correct extension for the
source file, as the compiler will use it to determine in which language the
source code is written. Indeed, the Ymir language was made to be interoperable
with the C language, and therefore the compiler is also able to compile C source
files, which use the \textit{.c} extension.

If the source code does not respect the syntax or semantics of the Ymir
language, the compiler, upon invocation, will issue error messages. These
messages describe the origin of the errors, and are made to be readable by the
programmer to help them correct the problems.
\vfill%
      
\section{How this book is structured}
\label{sec:basics:structure}

This book is two books, and the colour of the page says which one is in hand.

The first part, on white pages, is a course. Its chapters each gather a set of
related notions, and each notion appears at the point where it is first needed,
so the chapters are meant to be read in order. Every one of them ends with
exercises and their solutions.

The second part, on light grey pages, is the specification. It runs through the
same material in the same order, chapter for chapter, but formally and
exhaustively. It is for consulting rather than reading: a question about a notion
is answered in the part II chapter that matches the part I chapter where the
notion was taught.

Source code appears in listings. A badge in the top right corner of a listing
names the language it is written in, and the colour of its background says the
same thing again.

\begin{description}
\item[\ymirbadgeymir] \fcolorbox{black!35}{background!40}{\phantom{MM}}~a
  \textbf{yellow} background, for source code written in the Ymir language. This
  is by far the most common listing of the book.
\item[\ymirbadgemyil] \fcolorbox{black!35}{teal!10}{\phantom{MM}}~a
  \textbf{teal} background, for M-YIL code, and
  \fcolorbox{black!35}{olive!10}{\phantom{MM}}~a \textbf{green} one, for L-YIL
  code. These two languages are not written by the programmer, they are
  generated by the compiler, and are used in this book to describe what a Ymir
  program actually does once compiled. They are presented in
  Section~\ref{sec:basics:yil_language}.
\item[\ymirbadgeshell] \fcolorbox{black!35}{gray!10}{\phantom{MM}}~a
  \textbf{gray} background, for a terminal session. The lines starting with a
  prompt of the form \tokennolst{alice@dev:\mtildeascii{}\$} are commands typed
  by the user, and the lines that follow them are the output of those commands.
\item[\ymirbadgec] source code written in the \textbf{C} language, on the same
  background as an Ymir listing. It appears only where a program has to be
  linked against code that is not written in Ymir.
\end{description}

Some listings show a program that is wrong on purpose, so that the reason the
compiler rejects it can be explained. Such a listing is Ymir like any other, but
its badge is a warning sign, \ymirbadgeerror. Two highlights then appear inside
it: \hb{red} marks the code at fault, and \hcb{green} marks a correction applied
to it.

\vfill%
\pagebreak%

\section{A first program}
\label{sec:basics:first_program}

A Ymir program is made of symbols. A symbol is a named entity declared in a
source file -- a function, a variable, a type -- and the name is what the rest of
the program uses to refer to it. Of the kinds of symbol, the one to know first is
the function, because every Ymir program contains at least one: the function
named \token{main}, which the system calls when the program starts.

The following listing is a complete program. It writes \textit{Hello World!} to
the console, and then stops.

\begin{lstlisting}[style=coloredverbatim, caption=Source file \textit{hello.yr}, label=lst:basics:hello_world]
use std::io;

fn main() {
  println("Hello World!");
}
\end{lstlisting}

The first line declares that the program uses something from the \token{std::io}
module of the standard library -- here the function \token{println}. Standard
library modules are available without any further installation step, so a
\token{use} line is all it takes to reach one.

The keyword \token{fn} introduces a function, and \token{main} names it. Its
body, between the braces, is a single statement: a call to \token{println}. The
parentheses after a function name hold the arguments passed to it, in this case
one string literal. Every statement ends with a semicolon. Once the body runs out
of statements, \token{main} returns, and the program ends with it.

\token{println} finishes the text it writes with a line break, so that whatever
is written next starts on a new line. The standard library also provides
\token{print}, which writes exactly what it is given and nothing more. A line
break can equally be written inside the literal itself, as the escape sequence
\tokennolst{\textbackslash{}n}: escape sequences start with a backslash and stand
for characters that cannot be typed directly in a literal.
Chapter~\ref{chap:scalars} lists them. Note that a string literal is enclosed in
double quotes; single quotes are reserved for character literals, described in
that same chapter.

\subsection{Compilation and execution}
\label{sec:basics:compilation_execution}

A source file is not yet a program. The compiler, \token{gyc}, turns it into one:
given \textit{hello.yr}, the command below produces an executable file named
\textit{hello}.

\begin{lstlisting}[style=bashVerb, escapechar=@+]
@+\textcolor{teal!80}{alice@dev:\mtilde}@+$ gyc hello.yr -o hello
\end{lstlisting}

That file now sits in the current directory, and running it prints the expected
line.

\begin{lstlisting}[style=bashVerb, escapechar=@+]
@+\textcolor{teal!80}{alice@dev:\mtilde}@+$ ./hello
Hello World!
\end{lstlisting}


\vfill
\pagebreak

\section{Structure of a simple Ymir program}
\label{sec:basics:program_structure}

Most programs need more than one function. The next listing declares three --
\token{foo}, \token{bar} and \token{main} -- and calls them in a chain. Their
order in the file does not matter to the compiler: a function may call one
declared below it. Custom is nevertheless to put \token{main} last, or in a
module of its own.

\begin{lstlisting}[style=coloredverbatim, caption=First structured source file, label=lst:basics:fst_struct_src]
use std::io;

fn foo() {
  println("Called foo");
  bar(); // calling the bar function
}

fn bar() {
  // because we use std::io, println can be called directly
  println("Called bar");
}

/**
 * The entry point of the program
 */
fn main() {
  println("Program started");
  foo();
}
\end{lstlisting}

The listing also introduces comments. The compiler ignores them entirely; they
exist for whoever reads the code next. A later chapter shows how they can be
turned into documentation automatically. Two syntaxes are available:

\begin{itemize}
\item \token{//} comments out everything to the end of the line.
\item \token{/*} and \token{*/} enclose a comment that may span several lines.
  Starting each of those lines with a star, as on line~14 of
  Listing~\ref{lst:basics:fst_struct_src}, is a convention rather than a rule.
\end{itemize}

\subsection{Source code layout}
\label{sec:basics:source_layout}

The compiler reads a source file as a stream of tokens -- names, parentheses,
commas, brackets and the rest -- and treats any run of whitespace between two of
them as a single separator. Spaces, tabulations and line breaks are
interchangeable, and there may be any number of them. A function can therefore be
written on one line, or spread over a dozen, without changing its meaning.

%% restyle: skip -- spacing here is deliberately arbitrary, that is the point
\begin{lstlisting}[style=coloredverbatim, caption=Arbitrary code layout example]
use std::io;
fn foo(){println("Called foo");bar();}

fn
main
(   )
{
println ("Program started");
foo();
}
\end{lstlisting}

\subsection{Preferred layout}
\label{sec:basics:layout}

That freedom is worth spending on readability rather than on novelty, so the book
follows one layout throughout:

\begin{itemize}
\item one statement per line;
\item two spaces of indentation per level of scope, a scope being what the tokens
  \token{\{\}} delimit;
\item spaces around binary operators;
\item no space before \token{()} or \token{[]}, which attach to the name in front
  of them, as in \token{foo(a, b, c)};
\item no space before the colon that introduces a type, one space after it, as in
  \token{fn foo(a: i32)-> i32}.
\end{itemize}

Every listing in this book is written this way, and Ymir programs are best
written the same way -- uniform code is code a reader can skim.

\vfill
\pagebreak

\section{YIL language}
\label{sec:(chap1):yil_language}

Ymir is a high level programming language, meaning that a statement in Ymir can
represent multiple low level instructions. The C programming language is closer
to a machine level representation (even if there are still some high level
instructions that can be found in that language). In this book, in order to
fully understand the behavior of the Ymir language when executed on a machine,
and the memory representation a source code is describing, low level source code
close to the C language will be used. This low level source code uses a language
named YIL (Ymir Intermediate Language), and is a direct representation of the
high level source code in a low level form. A middle level representation will
also be used. In this middle level, some higher level operations are still
present, such as loops, or some high level memory management. In general it
won't be necessary to distinguish between the two levels of representation, as
they are rarely needed together, however if required, the terms L-YIL and
M-YIL, for respectively \textit{low} and \textit{middle} YIL, will be used.


As an example, the \textit{hello.yr} source file presented in
Listing~\ref{lst:hello_world} is transformed by the compiler into the following
M-YIL and L-YIL representations. The contents of these representations are not
yet discussed, but will be presented in the course of this document. The YIL
listings of this book are simplified for readability: they follow the structure
of what the compiler actually prints, but omit some of the internal annotations
it emits.

\begin{lstlisting}[style=myilVerb, caption=M-YIL representation of \textit{hello.yr}]
frame : [weak] std::io::println!{[c8]}::println (a(#8h) : [c8])-> void {
    std::io::print (a(#8h));
    std::io::putchar ('\n'c8);
    <unit-value>
}
frame : hello::main ()-> void {
    std::io::println!{[c8]}::println ("Hello World!"s8);
    <unit-value>
}
\end{lstlisting}

\begin{lstlisting}[style=lyilVerb, caption=L-YIL representation of \textit{hello.yr}]
frame :  [weak] _Y3std2io7printlnNS2c8FS2c8Zv (let a(#1) : (len-> u64, ptr-> *(u8)))-> void {
    _Y3std2io5printFS2c8Zv(a(#1));
    putchar(10);
}
frame :  _Y5hello4mainFZv ()-> void {
    _Y3std2io7printlnNS2c8FS2c8Zv((len-> 12, ptr-> "Hello World!"));
}
frame :  main (let argc(#1) : u32, let argv(#2) : *(void))-> i32 {
    _yrt_run_main(argc(#1), argv(#2), &_Y5hello4mainFZv);
    return 0;
}
\end{lstlisting}

YIL representations are automatically generated by the compiler, and are not
intended to be used as a programming language (in fact, there is no compiler to
convert their plain text representation into an executable). The
Ymir \textit{gyc} compiler generates the YIL representation when
the \textit{-fdump-ymir} option is specified. This option generates a set of
intermediate representations, among which two files with the
extensions \textit{.ydump-sem} and \textit{.ydump-yil} can be found, containing
the M-YIL and L-YIL representations of the compiled source file.

\begin{lstlisting}[style=bashVerb, escapechar=@+]
@+\textcolor{teal!80}{alice@dev:\mtilde}@+$ gyc -fdump-ymir hello.yr -o hello
@+\textcolor{teal!80}{alice@dev:\mtilde}@+$ ls
@+\textcolor{green!70!black}{hello}@+
hello.yr
hello.yr.ydump-decls.1
hello.yr.ydump-decls.2
hello.yr.ydump-opt
hello.yr.ydump-sem
hello.yr.ydump-yil
\end{lstlisting}

\vfill
\pagebreak

\section{Gyllir}
\label{sec:basics:gyllir}

Calling \token{gyc} by hand stops being reasonable at about the second source
file. Gyllir is the build system and package manager that takes over from there:
it assembles the command lines, tracks the version of the project, and fetches
the libraries it depends on. Nothing in Ymir requires it, and this book uses
whichever of the two tools is the more convenient -- \token{gyc} directly for the
short single-file listings that make up most of the examples, and whenever the
compiler's intermediate representations are on show, as in the previous section;
Gyllir as soon as a program spans several files or reaches for an external
library.

A new project starts with the \textit{init} subcommand, which asks a handful of
questions, each with a default, and writes out an empty project around the
answers.

\begin{lstlisting}[style=bashVerb, escapechar=@+]
@+\textcolor{teal!80}{alice@dev:\mtilde}@+$ mkdir my-project
@+\textcolor{teal!80}{alice@dev:\mtilde}@+$ cd my-project
@+\textcolor{teal!80}{alice@dev:\mtilde}@+@+\textcolor{green!60!black}{/my-project}@+$ gyllir init
Name [main]: my-project
Author name [Alice Doe <alice-doe@mail.com>]:
Description [A minimal Ymir app]:
License [proprietary]:
Type (executable/library) [executable]:
Registry [local:/home/alice/.local/gyllir/my-project]:
\end{lstlisting}

That leaves three files in the current directory.
\textit{gyllir.toml} is the configuration Gyllir reads on every invocation; with
the defaults above it describes an executable. \textit{src} holds the source
files -- for now a single \textit{main.yr} with placeholder content -- and
\textit{test} holds the unit tests, the subject of a later chapter.

\begin{lstlisting}[style=bashVerb, escapechar=@+]
@+\textcolor{teal!80}{alice@dev:\mtilde}@+@+\textcolor{green!60!black}{/my-project}@+$ tree
@+\color{teal}\hspace{0.5ex}\raisebox{0.5ex}{\texttt{.}}\vspace{-3pt}@+
@+\japan{├──} \raisebox{0.5ex}{\texttt{gyllir.toml}}\vspace{-3.5pt}@+
@+\japan{├──}~\color{teal}{\raisebox{0.5ex}{\texttt{src}}}\vspace{-3.5pt}@+
@+\japan{│~~~~~└──}~\raisebox{0.5ex}{\texttt{main.yr}}\vspace{-3.5pt}@+
@+\japan{└──}~\color{teal}{\raisebox{0.5ex}{\texttt{test}}}\vspace{-3.5pt}@+
@+\japan{~~~~~~~└──}~\raisebox{0.5ex}{\texttt{\_\_test\_\_.yr}}\vspace{-3.5pt}@+

2 directories, 3 files
\end{lstlisting}

\textit{gyllir build} produces the executable; \textit{gyllir run} produces it
and then runs it.

\begin{lstlisting}[style=bashVerb, escapechar=@+]
@+\textcolor{teal!80}{alice@dev:\mtilde}@+@+\textcolor{green!60!black}{/my-project}@+$ gyllir run
  @+\color{black!30!applegreen}{Compiling}@+ my-project v0.1.0
  @+\color{black!30!applegreen}{Compiling}@+ 1 modules in thread 1 (among 1 changed files of 1)
  @+\color{black!30!applegreen}{Produced}@+ executable file my-project
  @+\color{black!30!applegreen}{Executing}@+ project my-project v0.1.0

Hello World!
\end{lstlisting}

\vfill
\pagebreak

\section*{Exercises}
\markright{\MakeUppercase{Exercises}}%
\subsection*{Exercise 1}

Write a Ymir program that prints the following to the console:
\begin{verbatim}
Hello!
I wrote my first Ymir program.
Ain't that fun?
\end{verbatim}

\subsection*{Exercise 2}
The following program contains several errors:

\begin{lstlisting}[style=coloredverbatimError, escapechar=@, caption=Ymir program with errors]
*/ Comments of the main function @/*@
fn main {
  println("If this text")
  println("appears on the screen, ');
  println('you got it right').
)
\end{lstlisting}
Correct the errors and compile the program with \token{gyc} to test your
corrections.

\subsection*{Exercise 3}

What does the following program output to the console?
\begin{lstlisting}[style=coloredverbatim]
use std::io;

fn pause() {
  print("BREAK");
}

fn main() {
  println("\nDear reader, ");
  print("have a ");
  pause();
  println("!");
}
\end{lstlisting}

\subsection*{Exercise 4}

The following program is correct, and prints \textit{1 + 2 = 3} when it is
executed, but it does not follow the layout presented in
Section~\ref{sec:basics:layout}. Rewrite it so that it does.

%% restyle: skip -- the point of the exercise is that this layout is wrong
\begin{lstlisting}[style=coloredverbatim]
use std::io;
fn add (a : i32,b : i32)-> i32 {a+b}
fn main
(){
let x=add (1 , 2);
println ("1 + 2 = ",x);
}
\end{lstlisting}

\subsection*{Exercise 5}

The compiler was invoked on a source file named \textit{calc.yr}, and produced
the following representation. In which file was this representation written, and
which command line produced it? What does the program print when it is executed?

\begin{lstlisting}[style=myilVerb]
frame : calc::twice (a(#2) : i32)-> i32 return (a(#2) * 2)
frame : calc::main ()-> void {
    std::io::println!{[c8], i32}::println ("= "s8, calc::twice (21));
    <unit-value>
}
\end{lstlisting}

\subsection*{Exercise 6}

Using Gyllir, create a project named \textit{greeter} that prints
\textit{Hello reader!} when it is run. Which commands must be typed, and which
file of the project must be modified?

\vfill%
\pagebreak

\forceevenpage{}

\section*{Solutions}
\markright{\MakeUppercase{Solutions}}%
\subsection*{Exercise 1}

\begin{lstlisting}[style=coloredverbatim, caption=Solution for exercise 1]
use std::io;

fn main() {
  println("Hello!");
  println("I wrote my first Ymir program.");
  println("Ain't that fun?");
}
\end{lstlisting}

\subsection*{Exercise 2}

%% check: skip -- answer listing is written almost entirely in LaTeX escapes
\begin{lstlisting}[style=coloredverbatim, escapechar=@]
@\hcb{\color{purple}{use} \color{black}{std::io;}}@

@\hcb{/*}@ @\textit{Comments of the main function}@ @\hcb{*/}@
fn main@\hcb{()}@ {
  println(@\color{black!30!applegreen}{"If this text"}@)@\hcb{;}@
  println(@\color{black!30!applegreen}{"appears on the screen,}@ @\hcb{"}@);
  println(@\hcb{"}@@\color{black!30!applegreen}{you got it right}@@\hcb{"}@);
@\hcb{\}}@
\end{lstlisting}

\subsection*{Exercise 3}
The output starts with a line break.
\begin{lstlisting}[style=bashVerb]

Dear reader,
have a BREAK!
\end{lstlisting}

\subsection*{Exercise 4}

Each statement is written on its own line, the body of the two functions is
indented by two spaces, the parameter lists and the type colons are attached to
the name that precedes them, and the binary operators are surrounded by spaces.

\begin{lstlisting}[style=coloredverbatim, caption=Solution for exercise 4]
use std::io;

fn add(a: i32, b: i32)-> i32 {
  a + b
}

fn main() {
  let x = add(1, 2);
  println("1 + 2 = ", x);
}
\end{lstlisting}

\subsection*{Exercise 5}

This is a M-YIL representation, so it was written in the file
\textit{calc.yr.ydump-sem}, by the command line below. The L-YIL
representation of the same source file was written in \textit{calc.yr.ydump-yil}
at the same time.

\begin{lstlisting}[style=bashVerb, escapechar=@+]
@+\textcolor{teal!80}{alice@dev:\mtilde}@+$ gyc -fdump-ymir calc.yr -o calc
\end{lstlisting}

The \token{main} function calls \token{println} with two values, the string
\token{"= "} and the result of the call to \token{twice} with the value
\token{21}. The frame of \token{twice} returns its parameter multiplied by two,
so the program prints \textit{= 42}.

\subsection*{Exercise 6}

The \textit{init} subcommand creates the project, whose name is asked by the
first question. The file to modify is \textit{src/main.yr}, the source file
created by Gyllir, which prints \textit{Hello World!} until it is changed.

\begin{lstlisting}[style=bashVerb, escapechar=@+]
@+\textcolor{teal!80}{alice@dev:\mtilde}@+$ mkdir greeter
@+\textcolor{teal!80}{alice@dev:\mtilde}@+$ cd greeter
@+\textcolor{teal!80}{alice@dev:\mtilde}@+@+\textcolor{green!60!black}{/greeter}@+$ gyllir init
Name [main]: greeter
@+\textcolor{black!50}{... the other questions can keep their default value}@+
\end{lstlisting}

\begin{lstlisting}[style=coloredverbatim, caption=Content of \textit{src/main.yr}]
use std::io;

fn main() {
  println("Hello reader!");
}
\end{lstlisting}

The project is then compiled and executed by the \textit{run} subcommand.

\begin{lstlisting}[style=bashVerb, escapechar=@+]
@+\textcolor{teal!80}{alice@dev:\mtilde}@+@+\textcolor{green!60!black}{/greeter}@+$ gyllir run
  @+\color{black!30!applegreen}{Compiling}@+ greeter v0.1.0
  @+\color{black!30!applegreen}{Compiling}@+ 1 modules in thread 1 (among 1 changed files of 1)
  @+\color{black!30!applegreen}{Produced}@+ executable file greeter
  @+\color{black!30!applegreen}{Executing}@+ project greeter v0.1.0

Hello reader!
\end{lstlisting}

